%%%%%%% Resumen
\addchap{Resumen}

% Los motores de corriente continua (DC) son utlizados ampliamente en aplicaciones industriales y académicas por su sencillez de control y buena relación costo–prestaciones. La obtención precisa de sus parámetros eléctricos y mecánicos es esencial para el modelado, el diseño de controladores y el diagnóstico de fallas; sin embargo, los procedimientos tradicionales suelen ser manuales, lentos y propensos a errores. Este trabajo de investigación presenta el diseño e implementación de una máquina semiautomatizada para la estimación de parámetros de motores DC de hasta 12\,V, a partir de la excitación por voltaje de entrada y la medición de la velocidad angular como variable de salida.

El presente trabajo de investigación aborda el diseño e implementación X. En el primer capítulo se presentan los antecedentes, la motivación, los objetivos y el alcance de la máquina a diseñar. Los capítulos segundo y tercero desarrollan el marco teórico, explicando los conceptos fundamentales para el diseño de X. Además, se recopilan investigaciones previas relevantes y se comparan tecnologías utilizadas de X.

% El sistema propuesto estima de manera integrada la resistencia e inductancia del devanado, la constante de fuerza contraelectromotriz (voltaje–velocidad), la constante de torque (torque–corriente), el coeficiente de fricción viscosa y la inercia del eje. La solución se estructura como un banco de prueba compacto que genera señales de excitación, adquiere datos, ejecuta algoritmos de identificación y reporta los resultados con trazabilidad. Para ello, se desarrolla el diseño conceptual mediante la definición de requisitos, la estructura de funciones y la matriz morfológica; posteriormente, se detallan y materializan los dominios mecánico, electrónico, de software y de control del sistema.

% La validación incluye simulaciones y ensayos experimentales con motores DC comerciales dentro del rango de 12\,V, comparando los parámetros estimados frente a mediciones de referencia obtenidas mediante procedimientos de laboratorio. Los resultados evidencian repetibilidad y concordancia dentro de márgenes aceptables para propósitos de modelado y control, demostrando la viabilidad de una estimación semiautomatizada, rápida y segura. Finalmente, se discuten limitaciones (suposición de fricción predominantemente viscosa, cargas pasivas y régimen SISO) y se proponen extensiones futuras, como la incorporación de medición de corriente en línea, modelos de fricción no lineal y pruebas bajo variación de carga.

\clearpage
%%%%%%% Dedicatoria
\addchap{Dedicatoria}
\begin{quote}
Mi documento de tesis está dedicado a \lipsum[1]
\end{quote}
%%%%%%% Agradecimientos
\addchap{Agradecimientos}
\begin{quote}
Quisiera agradecer a todas las personas que me apoyaron durante el desarrollo de esta tesis. Especialmente gracias a (i) toda mi familia que me motivó a terminar este trabajo, (ii) a mi asesor de tesis xxxx, siempre estuvo disponible para guiarme en la dirección correcta, (iii) al responsable del Laboratorio yyyy, el Ingeniero Z, que me apoyo en el buen uso del equipo, y (iv) a ....
\end{quote}